\section{APPENDIX A}

\subsection{Paragraf 1: Teori Penetapan Harga}
\begin{frame}{Motivasi HJM}
    \begin{itemize}
        \item \textbf{1. Tantangan:} Kompleksitas pemodelan harga wajar instrumen derivatif.
        \item \textbf{2. Black-Scholes:} Titik acuan historis dengan limitasi fitur yang sangat terbatas.
        \item \textbf{3. Kegagalan Tradisional:} Ketidakmampuan menangkap seluruh faktor risiko pasar heterogen.
        \item \textbf{4. Solusi HJM:} Paradigma modern yang lebih representatif terhadap faktor risiko.
    \end{itemize}
\end{frame}

\begin{frame}{Pertanyaan yang dijawab pada Paragraf 1}
    \centering
    \small
    \begin{tabular}{lp{4.5cm}p{6cm}}
        \toprule
        Kalimat ke & Pertanyaan & Jawaban \\
        \midrule
        1 & Harga derivatif sulit? & Tantangan pemodelan matematis sangat kompleks. \\
        3 & Kelemahan model tradisional? & Gagal menangkap faktor risiko finansial secara akurat. \\
        4 & Mengapa HJM diminati? & Atasi limitasi model tradisional di skenario rumit. \\
        \bottomrule
    \end{tabular}
\end{frame}

\subsection{Paragraf 2: Optimasi Monte Carlo}
\begin{frame}{Tahapan Desain}
    \begin{itemize}
        \item \textbf{1. Komputasi:} Beban berat penyelesaian model multifaktor.
        \item \textbf{2. Monte Carlo:} Metode dominan namun inefisien (\textit{low convergence}).
        \item \textbf{3. Dua Tahap:} Reduksi kompleksitas (PCA) diikuti percepatan konvergensi.
        \item \textbf{4. Alternatif QC:} Teknologi kunci akselerasi kedua tahap optimasi.
    \end{itemize}
\end{frame}

\begin{frame}{Pertanyaan yang dijawab pada Paragraf 2}
    \centering
    \small
    \begin{tabular}{lp{4.5cm}p{6cm}}
        \toprule
        Kalimat ke & Pertanyaan & Jawaban \\
        \midrule
        1 & Beban komputasi multifaktor? & Sangat berat, perlu sumber daya tangguh. \\
        3 & Desain simulasi MC efisien? & Reduksi kompleksitas (PCA) dan akselerasi konvergensi. \\
        4 & Peran QC di sini? & Akselerasi kedua tahap optimasi tersebut. \\
        \bottomrule
    \end{tabular}
\end{frame}

\subsection{Paragraf 3: Elemen Fundamental P(t,T)}
\begin{frame}{Variabel Kunci Obligasi}
    \begin{itemize}
        \item \textbf{P(t,T):} Fungsi harga obligasi pada waktu t sebagai variabel fundamental.
        \item \textbf{P(T,T) = 1:} Kondisi batas saat jatuh tempo untuk konsistensi model.
        \item \textbf{Risk-Neutral:} Menghubungkan kurva harga dengan prinsip \textit{risk-neutral pricing}.
        \item \textbf{Faktor Diskon:} Valuasi arus kas masa depan dan inversnya sebagai faktor kapitalisasi.
    \end{itemize}
\end{frame}

\begin{frame}{Pertanyaan yang dijawab pada Paragraf 3}
    \centering
    \small
    \begin{tabular}{lp{4.5cm}p{6cm}}
        \toprule
        Kalimat ke & Pertanyaan & Jawaban \\
        \midrule
        1 & Definisi obligasi nol kupon? & Kontrak jaminan bayar 1 unit mata uang pada T. \\
        2 & Nilai saat jatuh tempo? & Harus sama dengan satu ($P(T,T)=1$). \\
        4 & Interpretasi ekonomi? & Sebagai kurva faktor diskon arus kas masa depan. \\
        \bottomrule
    \end{tabular}
\end{frame}

\subsection{Paragraf 4: Suku Bunga Sesaat (Short Rate)}
\begin{frame}{Definisi Suku Bunga}
    \begin{itemize}
        \item \textbf{r(t):} Suku bunga sesaat atau \textit{instantaneous spot rate}.
        \item \textbf{B(t):} Persamaan akun pasar uang sebagai instrumen akumulasi tanpa risiko.
        \item \textbf{Stokastik:} Suku bunga masa depan dimodelkan melalui proses stokastik.
    \end{itemize}
\end{frame}

\begin{frame}{Pertanyaan yang dijawab pada Paragraf 4}
    \centering
    \small
    \begin{tabular}{lp{4.5cm}p{6cm}}
        \toprule
        Kalimat ke & Pertanyaan & Jawaban \\
        \midrule
        1 & Apa itu suku bunga sesaat? & Pengembalian bebas risiko pada periode infinitesimal. \\
        3 & Akumulasi pasar uang? & Fungsi eksponensial integral suku bunga sesaat. \\
        \bottomrule
    \end{tabular}
\end{frame}

\subsection{Paragraf 5: Kepastian vs Ketidakpastian}
\begin{frame}{Sifat Deterministik dan Acak}
    \begin{itemize}
        \item \textbf{Waktu t:} Pengamatan pasar saat ini bersifat deterministik.
        \item \textbf{Waktu T:} Butuh pemodelan stokastik untuk prediksi masa depan.
        \item \textbf{Risk-Neutral Framework:} Formula penetapan harga obligasi berbasis $B(t)$.
        \item \textbf{Random Variables:} Variabel acak pada jatuh tempo sebagai sumber risiko harga.
    \end{itemize}
\end{frame}

\begin{frame}{Pertanyaan yang dijawab pada Paragraf 5}
    \centering
    \small
    \begin{tabular}{lp{4.5cm}p{6cm}}
        \toprule
        Kalimat ke & Pertanyaan & Jawaban \\
        \midrule
        1 & Kepastian suku bunga di t? & Deterministik, diobservasi langsung di pasar. \\
        3 & Hubungan obligasi \& B(t)? & Melalui ekspektasi ukuran martingal ekuivalen $Q_B$. \\
        5 & Lokasi ketidakpastian? & Pada nilai suku bunga masa depan (variabel acak). \\
        \bottomrule
    \end{tabular}
\end{frame}

\subsection{Paragraf 6-10: Pemodelan HJM Tingkat Lanjut}
\begin{frame}{Dinamika dan Faktor Noisy}
    \begin{itemize}
        \item \textbf{Klasifikasi Model (P6):} Berdasarkan jumlah faktor gangguan (\textit{noise factors}).
        \item \textbf{HJM Solution (P7):} Solusi sistematis pemodelan \textit{forward rates} secara langsung.
        \item \textbf{Matematika HJM (P8):} Hubungan turunan parsial antara $f(t,T)$ dan $P(t,T)$.
        \item \textbf{No-Arbitrage (P9):} Koefisien \textit{drift} $\alpha$ menjamin kondisi bebas arbitrase.
        \item \textbf{PCA (P10):} Reduksi dimensi untuk mengatasi biaya komputasi simulasi numerik.
    \end{itemize}
\end{frame}

\begin{frame}{Pertanyaan yang dijawab pada Paragraf 6-10}
    \centering
    \small
    \begin{tabular}{lp{4.5cm}p{6cm}}
        \toprule
        Kalimat ke & Pertanyaan & Jawaban \\
        \midrule
        P6.3 & Mengapa satu faktor ditinggal? & Gagal representasi distribusi gabungan antar maturitas. \\
        P7.2 & Solusi HJM? & Pemodelan suku bunga masa depan (\textit{forward}) langsung. \\
        P9.1 & Fungsi \textit{drift} $\alpha$? & Menjamin kondisi bebas arbitrase dalam model. \\
        P10.2 & Cara reduksi dimensi? & PCA untuk dapat \textit{eigenvector} utama yang mewakili dinamika. \\
        \bottomrule
    \end{tabular}
\end{frame}
